\documentclass[10pt,a4paper]{amsart}
\usepackage[margin=19mm]{geometry}
\usepackage{amsmath,amssymb,mathtools}
\usepackage[scr=rsfs,cal=euler]{mathalfa}
\usepackage{xfrac}
\usepackage[unicode,hidelinks,bookmarks=false]{hyperref}
\newcommand{\ie}{i.e.\ }
\newcommand{\cf}{cf.\ }
\newcommand{\resp}{respectively\ }
\newcommand{\Subsection}{\subsection}
\providecommand{\nobreakdash}{\nobreak\hbox{-}}
\setlength{\emergencystretch}{2em}
\multlinegap0pt
\usepackage{amssymb}
\usepackage{mathtools}
\usepackage{enumerate}
\newtheorem{theorem}{Theorem}
\newcommand{\id}{\operatorname{id}}
\title{Weyl structures for path geometries}
\author{Andreas \v Cap and Zhangwen Guo}
\date{Source: arXiv:2604.12846v1 (CC BY 4.0)}
\begin{document}
\maketitle

\noindent\emph{Source and licence.} Weyl structures for path geometries. Version arXiv:2604.12846v1 (CC BY 4.0), distributed by arXiv under the Creative Commons Attribution 4.0 International licence, http://creativecommons.org/licenses/by/4.0/, as recorded in the arXiv licence metadata for this exact version and shown on the abstract page https://arxiv.org/abs/2604.12846v1. Full licence notice: \texttt{licenses/arxiv-2604.12846-CC-BY-4.0.txt}.\par\smallskip

\noindent\textit{Editorial scope.} Counted unit: the article's theorem tors-curv (source0.tex lines 304--315). The article's complete original proof of it is delivered with it (source0.tex lines 319--345, one \texttt{proof} environment of 27 lines, which the article places away from the statement). No mathematics is written, repaired or re-derived: the statement and the proof are the article's own lines, in the article's own notation, and no retained line is substituted or re-flowed.  Retained uncounted as the source-local context the proof needs: lines 222--232; lines 248--258; lines 276--281.  Nothing was omitted from any retained span, so the package carries no omission record.  No retained line contains a citation, so the wrapper carries no bibliography and no bibliography entry.  No retained line contains a figure environment, an image-inclusion command or drawing code, so no image is delivered and the package declares no asset dependency.  Editorial formatting only, in addition: an AMS \texttt{amsart} preamble that restates the article's own declarations and macros used by the retained lines, namely the environments \texttt{theorem} on separate counters, together with the standard packages those lines need.  The retained environments print in this document as Theorem 1 in order of appearance, whereas the article prints the same items with the numbering of its own full document.  The complete native source of the article as distributed in the arXiv e-print for this exact version is bundled unchanged as \texttt{sources/198430/source0.tex}.
\begin{theorem}\label{thm:main}
    Let $(M,E,V)$ be a path geometry and let $\xi_0$ be a nowhere vanishing local smooth section of the subbundle $E\subset TM$ defined on an open subset $U\subset M$. Then there are a unique linear connection $\nabla$ on $TU$ and a unique decomposition $TU=E\oplus V\oplus \iota(Q)$ such that the following conditions are satisfied: 
    \begin{itemize}
        \item[(i)] The subbundles $E$, $V$ and $\iota(Q)$ of $TU$ are parallel for $\nabla$, and $\nabla\xi_0=0$. 
        \item[(ii)] The bundle map $\iota\circ\mathcal L:E\otimes V\to\iota(Q)$ is parallel for (the connections induced by) $\nabla$.
        \item[(iii)] The torsion $\tau$ of $\nabla$ satisfies $\tau(E,V)\subset \iota(Q)$ and $\tau(H,\iota(Q))\subset H$, and for sections $\eta_1,\eta_2\in\Gamma(V)$ and putting $\zeta_i:=\iota(\mathcal L(\xi_0,\eta_i))\in\Gamma(\iota(Q))$ we get 
        \begin{equation}\label{eq:norm}
            \mathcal L(\tau(\xi_0,\zeta_1),\eta_2)=2\mathcal L(\xi_0,\tau(\eta_1,\zeta_2)).
        \end{equation}
    \end{itemize}
\end{theorem}

\begin{equation}\label{homog1-props}
\begin{aligned}
\begin{cases}
\nabla^E\xi&=df\otimes\xi_0\\
\mathcal L(\xi_0,\nabla^V_{\xi}\eta)=\mathcal L(\xi_0,\Pi_V([\xi,\eta]))&=\frac 1 2 q([\xi,[\xi_0,\eta]])\\
\Pi_E([\xi,\eta])&=-df(\eta)\xi_0\\
\mathcal L(\xi_0,\nabla^V_{\eta_1}\eta_2)&=q([\eta_1,[\xi_0,\eta_2]])\\
\nabla^{\iota(Q)}\iota(\mathcal L(\xi,\eta))&=\iota(\mathcal L(\nabla^E\xi,\eta)+\mathcal L(\xi,\nabla^V\eta)).
\end{cases}
\end{aligned}
\end{equation}

    \begin{equation}\label{homog2-props}
    \begin{aligned}
     \nabla^V_{\zeta}\tilde\eta&=\Pi_V([\zeta,\tilde\eta])+\tfrac12\alpha([\xi_0,[\xi_0,\tilde\eta]])\eta\\
     \nabla^{\iota(Q)}\zeta &=\iota(\mathcal L(\xi_0,\nabla^V\eta)).    
    \end{aligned}
    \end{equation}

\begin{theorem}\label{tors-curv}
  Let $(M,E,V)$ be a path geometry, $\xi_0\in\Gamma(E)$ a nowhere vanishing section, $(\nabla,\Pi)$ the data associated to $\xi_0$ in Theorem \ref{thm:main} and let $\tau$ and $R$ be the torsion and the curvature of $\nabla$. Then we have $\tau(V,V)=0$ and $\tau(TM,\iota(Q))\subset H$. Moreover, denoting by $\alpha\in\Omega^1(M)$ the form charaterized by $\Pi_E=\alpha\xi_0$, we get for $\eta_1,\eta_2,\eta_3\in\Gamma(V)$ and $\zeta_i=\iota(\mathcal L(\xi_0,\eta_i))$ for $i=1,2,3$ the following:
\begin{equation}\label{eq:tors-curv}
\begin{aligned}
\begin{cases}
\tau(\eta_1,\zeta_2)&=-\tfrac12\alpha([\xi_0,\zeta_1])\eta_2-\alpha([\eta_1,\zeta_2])\xi_0\\ 
R(\xi_0,\eta_1)(\eta_2)&=-\tfrac12\alpha([\xi_0,\zeta_1])\,\eta_2-\alpha([\xi_0,\zeta_2])\,\eta_1\\
R(\eta_1,\eta_2)(\eta_3)&=-\alpha([\eta_1,\zeta_3])\,\eta_2+\alpha([\eta_2,\zeta_3])\,\eta_1.
\end{cases}
\end{aligned}
\end{equation}  
\end{theorem}

\begin{proof}[Proof of Theorem \ref{tors-curv}]
    By the Jacobi identity $0=[\eta_1,[\xi_0,\eta_2]]-[\eta_2,[\xi_0,\eta_1]]-[\xi_0,[\eta_1,\eta_2]]$. By involutivity, we get $[\eta_1,\eta_2]\in\Gamma(V)$ so applying $q$ to this equation and using the fourth line of \eqref{homog1-props}, we get $\mathcal L(\xi_0,\tau(\eta_1,\eta_2))=0$. Since $\tau(V,V)\subset V$ by definition, this implies $\tau(V,V)=0$. 
    
    Since we know that $\tau(\eta_1,\zeta_2)=\nabla_{\eta_1}\zeta_2-\nabla_{\zeta_2}\eta_1-[\eta_1,\zeta_2]$ lies in $\Gamma(H)$, it has to be given by $-\nabla_{\zeta_2}\eta_1-\Pi_V([\eta_1,\zeta_2])-\Pi_E([\eta_1,\zeta_2])$. By the first line of \eqref{homog2-props}, the first two summands give $-\tfrac12\alpha([\xi_0,[\xi_0,\eta_1]])\eta_2$. As we have observed already we can replace $[\xi_0,\eta_1]$ with $\zeta_1$ without changing the result, so the first line of \eqref{eq:tors-curv} follows readily. 

    The fourth line of \eqref{homog1-props} shows that  
    $\mathcal L(\xi_0,\nabla_{\eta_1}\nabla_{\eta_2}\eta_3)=q([\eta_1,[\xi_0,\nabla_{\eta_2}\eta_3]])$.
    It also shows that $q([\xi_0,\nabla_{\eta_2}\eta_3])=q([\eta_2,[\xi_0,\eta_3]])$, so $[\xi_0,\nabla_{\eta_2}\eta_3]-[\eta_2,[\xi_0,\eta_3]]+\Pi_E([\eta_2,[\xi_0,\eta_3]])\in\Gamma(V)$ by the third line of \eqref{homog1-props}. By involutivity of $V$, we conclude that
    $$
    \mathcal L(\xi_0,\nabla_{\eta_1}\nabla_{\eta_2}\eta_3)=q([\eta_1,[\eta_2,[\xi_0,\eta_3]]]) + \alpha([\eta_2,[\xi_0,\eta_3]])\mathcal L(\xi_0,\eta_1). 
    $$
    Subtracting the same term with $\eta_1$ and $\eta_2$ exchanged as well as $\mathcal L(\xi_0,\nabla_{[\eta_1,\eta_2]}\eta_3)=q([[\eta_1,\eta_2],[\xi_0,\eta_3]])$, the first terms in the right hand side cancel. Since $R(\eta_1,\eta_2)(\eta_3)$ is a section of $V$, this implies the last line of \eqref{eq:tors-curv}. 

    To deal with the second line of \eqref{eq:tors-curv}, we first observe that both sides of the claimed identity are linear over smooth functions in $\eta_1$ and $\eta_2$, so it suffices to verify them locally for the elements of a frame. Locally around each point, we can fix a hypersurface $\tilde M\subset M$ whose tangent spaces are transversal to $E$, then choose a local frame of $V$ along $\tilde M$ and extend it by parallel transport along $\xi_0$ to a local frame defined on an open subset. The frame elements then are parallel in $E$-directions, so we can prove our identity assuming that $\nabla_{\xi_0}\eta_i=0$ for $i=1,2$. This also implies that $\Pi_V([\xi_0,\eta_i])=0$ for $i=1,2$ and hence $\zeta_i=[\xi_0,\eta_i]$ and $\nabla_{\xi_0}\zeta_i=0$ for $i=1,2$. Hence the second line of \eqref{homog2-props} shows that 
    \begin{equation}\label{tech1}
    -\mathcal L(\xi_0,\nabla_{[\xi_0,\eta_1]}\eta_2)=-\mathcal L(\xi_0,\nabla_{\zeta_1}\eta_2)=-q([\xi_0,\Pi_V([\zeta_1,\eta_2]))-\tfrac12\alpha([\xi_0,[\xi_0,\eta_2]])\mathcal L(\xi_0,\eta_1). 
    \end{equation}
    Writing $\Pi_{-2}:=\id-\Pi$, we can write $\Pi_V([\zeta_1,\eta_2])=[\zeta_1,\eta_2]-\Pi_E([\zeta_1,\eta_2])-\Pi_{-2}([\zeta_1,\eta_2])$, and plugging into $q([\xi_0,\_])$ we can ignore the $\Pi_E$-term. Since $q(\tau(V,\iota(Q)))=0$, we get $\Pi_{-2}([\zeta_1,\eta_2])=-\nabla_{\eta_2}\zeta_1$ and since $q(\tau(E,\iota(Q)))=0$, we get $q([\xi_0,\nabla_{\eta_2}\zeta_1])=q(\nabla_{\xi_0}\nabla_{\eta_2}\zeta_1)=\mathcal L(\xi_0,\nabla_{\xi_0}\nabla_{\eta_2}\eta_1)$, where in the last step we have used that $\iota\circ \mathcal L(\xi_0,\_)$ is parallel. Hence we conclude that 
    \begin{equation}\label{tech2}
      -q([\xi_0,\Pi_V([\zeta_1,\eta_2]))=-q([\xi_0,[\zeta_1,\eta_2]])-\mathcal L(\xi_0,\nabla_{\xi_0}\nabla_{\eta_2}\eta_1). 
    \end{equation}
    To compute $\mathcal L(\xi_0,R(\xi_0,\eta_1)\eta_2)$, we have to add $\mathcal L(\xi_0,\nabla_{\xi_0}\nabla_{\eta_1}\eta_2)$ to \eqref{tech1}, and in view of $\tau(V,V)=0$, this adds up with the last term of \eqref{tech2} to $\mathcal L(\xi_0,\nabla_{\xi_0}[\eta_1,\eta_2])=\tfrac12q([\xi_0,[\xi_0,[\eta_1,\eta_2]]])$. Hence 
    \begin{equation}\label{tech3}
      \mathcal L(\xi_0,R(\xi_0,\eta_1)\eta_2)=\tfrac12q([\xi_0,[\xi_0,[\eta_1,\eta_2]]])-q([\xi_0,[\zeta_1,\eta_2]]) -\tfrac12\alpha([\xi_0,[\xi_0,\eta_2]])\mathcal L(\xi_0,\eta_1).
    \end{equation}
    The Jacobi identity and $\zeta_i=[\xi_0,\eta_i]$ readily show that the first two terms in the right hand side add up to $-\frac12 q([\xi_0,[\zeta_1,\eta_2]])-\frac12q([\xi_0,[\zeta_2,\eta_1]])=-\frac12 q([[\xi_0,\zeta_1],\eta_2])-\frac12q([[\xi_0,\zeta_2],\eta_1])$. Now in the defining equation for $\tau(\xi_0,\zeta_i)$ the derivative terms both vanish, so $-[\xi_0,\zeta_i]=\tau(\xi_0,\zeta_i)\in\Gamma(H)$. Forming the bracket with $\eta_j$ and applying $q$, only the $E$-component matters and this is given by $\alpha([\xi_0,\zeta_i])\xi_0$.  
\end{proof}    

\end{document}
